\chapter{100 câu cực viral kiểu đọc là muốn chụp đăng}
\markboth{100 câu cực viral kiểu đọc là muốn chụp đăng}{100 câu cực viral kiểu đọc là muốn chụp đăng}

\begin{truyen}{Lưu ý khi dùng 100 câu này}{How to use these lines}
\chuhoa{Đ}{ây là nhóm câu ngắn, gọn, đập thẳng và dễ lan truyền. Chúng hợp để đặt đầu chương, cuối bài đăng, story, caption, poster lớp, hoặc làm trích dẫn quảng bá cho sách.}
\textit{(These are short, sharp, shareable lines. They work well for chapter openers, post captions, stories, posters, or promotional quote cards.)}
\end{truyen}

\section{25 câu chụp lên Facebook là có người gật đầu ngay}
\begin{enumerate}[leftmargin=*]
\item Lớn lên không khó bằng nhận ra mình từng sống ảo với chính mình.
\item Người làm em mệt nhất nhiều khi là phiên bản dễ dãi của em.
\item Không phải em thiếu thời gian, em đang thiếu ưu tiên thật.
\item Cái tôi càng to, cuộc đời càng chật chỗ cho sự thật.
\item Có những ngày không cần giỏi hơn ai, chỉ cần đừng phản bội mình.
\item Đừng chờ đời nhẹ tay nếu em còn quá mềm với chính mình.
\item Bớt diễn ổn, sống thật sẽ đỡ mệt hơn nhiều.
\item Cảm hứng giúp bắt đầu, kỷ luật mới giữ em không gãy.
\item Có người rời đi không phải để phạt em, mà để dạy em.
\item Sĩ diện giữ mặt vài phút, năng lực giữ đời nhiều năm.
\item Đừng xin lỗi vì mình đang lớn lên theo hướng không còn vừa ý người khác.
\item Không ai cứu nổi người vẫn nghiện tự phá mình.
\item Đời không sợ em chậm, đời sợ em đứng yên mà cứ tưởng đang nghĩ lớn.
\item Càng thiếu nội lực càng thích sống bằng tiếng vỗ tay.
\item Mất ngủ vì ước mơ nghe hay, mất ngủ vì trì hoãn mới đúng đời hơn.
\item Đừng lấy cái sai của số đông làm thuốc tê cho lương tâm mình.
\item Người tử tế không phải người dễ dãi.
\item Chữa lành không phải trang trí nỗi đau cho đẹp hơn.
\item Cái giá của trưởng thành là bớt được vài ảo tưởng rất yêu thích.
\item Không phải bài học nào cũng đến bằng lời nhẹ nhàng.
\item Sống sâu không ồn, nhưng rất khó làm giả.
\item Đừng để vài phút cảm xúc phá vài năm đáng sống.
\item Trưởng thành là thôi đổ lỗi cho quá khứ trong mọi quyết định hiện tại.
\item Bản lĩnh là đau rồi vẫn không sống cộc đi.
\item Muốn đời nể mình, trước hết đừng sống kiểu chính mình cũng khó nể.
\end{enumerate}

\section{25 câu ngắn, cứng, dễ thành caption}
\begin{enumerate}[leftmargin=*,resume]
\item Sống cho rõ, đỡ phải chữa cháy cho những ngày mơ hồ.
\item Giữ bình yên khó hơn giữ sĩ diện.
\item Người mạnh là người không cần ồn để thấy mình có lực.
\item Đừng than cô đơn nếu em đang bỏ rơi chính mình trước nhất.
\item Mọi cú ngã đều hỏi một câu: em lớn lên hay cay đi.
\item Có chiều sâu rồi em sẽ bớt thích mấy thứ hào nhoáng ngắn hạn.
\item Đừng chê đường dài khi em còn sống nhờ đường tắt tưởng tượng.
\item Không ai nợ em một cuộc đời dễ hiểu.
\item Cái đáng sợ không phải là thất bại, mà là quen sống né sự thật.
\item Tỉnh táo là món quà đắt, nhưng rẻ hơn ảo tưởng rất nhiều.
\item Đừng dùng lòng tốt để bào mòn ranh giới.
\item Người đi xa hiếm khi đi bằng tâm trạng.
\item Có những mối quan hệ không xấu, chỉ là không còn đúng nữa.
\item Tự trọng không nuôi được cái tôi, nhưng cứu được nhân cách.
\item Đừng làm người lớn xác với trái tim chưa chịu trưởng thành.
\item Mỗi lần chọn điều đúng là một lần bớt sợ chính mình.
\item Không có gì ngầu trong việc làm khổ mình mà gọi đó là cá tính.
\item Im lặng đúng lúc cũng là một dạng bản lĩnh.
\item Đừng chờ mình hết sợ mới bắt đầu sống đàng hoàng.
\item Người thật sự tốt cho em sẽ không bắt em nhỏ lại để vừa họ.
\item Cần được yêu là thật, nhưng cần tự cứu còn thật hơn.
\item Học cách ở một mình trước khi đòi ở đúng với ai đó.
\item Đời không hứa công bằng, nên em càng phải rõ giá trị của mình.
\item Bớt giải thích, làm cho rõ hơn.
\item Bình yên thật không cần diễn.
\end{enumerate}

\section{25 câu kiểu đọc xong muốn lưu lại}
\begin{enumerate}[leftmargin=*,resume]
\item Có những ngày em không cần được hiểu hết, chỉ cần đừng tự bỏ mình.
\item Mình không yếu, mình chỉ từng quen sống thiếu điểm tựa bên trong.
\item Người làm được việc nhỏ đều mới gánh nổi việc lớn.
\item Đừng ghét bản thân trong giai đoạn mình đang học lớn.
\item Càng hiểu đời càng bớt nói câu nào cũng như chân lý.
\item Không phải nỗi buồn nào cũng cần kể, nhưng nỗi buồn nào cũng cần được hiểu.
\item Đôi khi thứ cần cắt không phải người khác, mà là thói quen tự phá mình.
\item Một người trưởng thành biết khi nào cần mềm và khi nào không được nhũn.
\item Người khôn không sống chỉ để chứng minh mình đúng.
\item Có những cánh cửa đóng lại để em thôi đi lạc.
\item Tương lai không cần em hứa nhiều, nó cần em làm đều.
\item Đừng vì vài lần thất vọng mà biến mình thành người mất hết khả năng tin.
\item Người biết giữ lời với bản thân thường bớt sợ đời hơn.
\item Chậm mà chắc đáng sợ hơn nhiều kiểu bùng nổ rồi tắt.
\item Điều đẹp nhất của một người trẻ là còn biết sửa mình.
\item Đừng dùng bận rộn để che nỗi sợ ngồi yên nhìn thẳng vào mình.
\item Có nỗi đau phải đi qua mới biết mình mạnh kiểu nào.
\item Không ai sống thay được đoạn em cần tự tỉnh.
\item Được công nhận là vui, tự công nhận mình xứng đáng còn quan trọng hơn.
\item Sống tử tế với mình là ngừng đưa mình vào những chỗ phải cầu xin sự tôn trọng.
\item Mỗi ngày bớt giả một chút là một ngày đáng kể.
\item Đừng lấy quá khứ làm nhà, lấy nó làm bản đồ thôi.
\item Người càng hiểu giá trị của bình yên càng ít thích những cuộc ồn ào vô nghĩa.
\item Không phải ai đến cũng để ở lại, nhưng ai cũng để lại gì đó cho em học.
\item Sống sâu là khi em không còn muốn thắng bằng mọi giá.
\end{enumerate}

\section{25 câu rất hợp làm trích dẫn quảng bá sách}
\begin{enumerate}[leftmargin=*,resume]
\item Cục súc không phải để làm đau, mà để đập vỡ lớp chống chế quá dày.
\item Có những câu thô mới đủ sức kéo người trẻ dừng lại mà nghĩ.
\item Một câu nói thật đôi khi cứu được cả tháng sống giả.
\item Điều khó nghe nhất thường là điều cần nghe nhất.
\item Sách này không vỗ về cái tôi, nó kéo cái tôi xuống đất.
\item Đừng đọc để thấy mình ngầu, đọc để thấy mình còn chỗ nào phải sửa.
\item Cái thấm không nằm ở câu chữ đao to búa lớn, mà ở chỗ nó chạm đúng nỗi mình né.
\item Có những câu đọc xong thấy bị xúc phạm, thật ra là bị chạm đúng.
\item Người trẻ cần sách không chỉ hay, mà phải đủ thật để không ngủ quên sau vài trang.
\item Một cuốn sách có ích là cuốn khiến em bớt biện hộ cho phần yếu của mình.
\item Viral thì dễ, ở lại trong đầu người đọc mới khó.
\item Câu hay không cứu được ai nếu nó chỉ dừng ở lượt chia sẻ.
\item Sự thật cứng nhưng đỡ hại hơn lời ngọt rỗng.
\item Đừng chỉ sưu tầm quote, hãy để vài câu sửa cách em sống.
\item Một câu đúng lúc đôi khi đáng hơn mười trang an ủi mờ.
\item Người trẻ không thiếu nội dung, họ thiếu nội dung dám nói thật.
\item Câu nào càng ngắn mà càng khiến người ta im vài giây thì càng có lực.
\item Thô vừa đủ, thật đủ sâu, đó là chỗ câu chữ bắt đầu bán được.
\item Sách chạm người trẻ phải nói trúng cái họ đang sống, không phải cái người lớn muốn họ nghe.
\item Một câu viral tốt là câu đọc lên thấy chính mình trong đó.
\item Đừng viết để được khen sâu, hãy viết để người ta thấy mình.
\item Câu chữ bán được không phải câu làm màu, mà là câu khiến người đọc muốn giữ lại.
\item Người ta chia sẻ thứ làm họ thấy bị hiểu hoặc bị đánh thức.
\item Muốn lan truyền, trước hết phải đủ thật.
\item Câu nào làm người đọc vừa đau vừa gật đầu, câu đó có đời sống riêng.
\end{enumerate}
